% !TeX root = fundamentals.tex
\section{Решение физической задачи}

Ответ на физическую задачу — это не только число. Нужно понять, почему
искомая величина связана с данными именно так, а затем показать эту связь
в решении. Для этого описание ситуации переводят на язык физических
величин и законов. Вычисления завершают рассуждение.

В этом параграфе договоримся об оформлении наших учебных решений:
сначала получаем выражение для искомой величины в буквенном виде,
затем подставляем числа. Разберём, что даёт такой порядок и как сделать
запись понятной другому человеку.

\subsection{От условия к уравнению}

Прежде чем записывать уравнения, полезно представить происходящее:
какое тело или систему мы рассматриваем, что с ними меняется и какие
условия остаются неизменными. Например, слова «температура постоянна»
задают связь между начальным и конечным состояниями. Это такая же
существенная часть условия, как численные значения величин.

Краткая запись \textit{«Дано»} собирает известные величины и условия,
а запись \textit{«Найти»} уточняет цель решения. Если требуется отношение
двух величин, именно его и следует указать: например, $n_2/n_1$, а не
два отдельных значения $n_1$ и $n_2$. Приводя данные к СИ, сохраняем
физический смысл величин: температура в уравнении состояния газа
должна быть абсолютной, а проценты при необходимости переводятся
в доли единицы.

Следующий шаг — выбор \textit{физической модели}. Уравнение относится
не просто к числам из условия, а к определённому телу или процессу.
Так, зависимость $s=vt$ описывает путь при постоянной скорости;
если скорость меняется, пользоваться этой записью с произвольным
значением $v$ уже нельзя.

Поэтому перед уравнением появляется короткое пояснение:
«Движение равномерное, поэтому путь равен произведению скорости
на время». Эта фраза объясняет и выбор зависимости, и возможность
её применения. Если в задаче важны направления сил, движения или
перемещения, их показывают на рисунке и выбирают оси.

\subsection{Зачем решать в общем виде}

Пусть известны путь $s$ и постоянная скорость $v$, а требуется найти
время движения $t$. Из зависимости $s=vt$ получаем
\[
t=\frac{s}{v}.
\]
Это и есть решение в общем виде: искомая величина выражена через
данные задачи. Теперь видно не только, какие действия выполнить
с числами, но и как устроена зависимость. При том же пути увеличение
скорости вдвое уменьшает время вдвое.

Буквы сохраняют смысл величин на протяжении решения.
Для двух состояний удобно использовать индексы: $T_1$, $T_2$;
для начальной скорости — $v_0$. Буква $x$ вполне уместна как
обозначение координаты. Важно, чтобы каждое обозначение было
определено и затем использовалось в том же смысле.

Если вместо $s$, $v$ и $t$ сразу записать числа и неизвестное $x$,
алгебра может оказаться верной, но физический смысл станет менее
видимым. Сам по себе выбор буквы не делает решение неправильным.
Проблема возникает, когда из записи нельзя понять, какая величина
найдена и на каком законе основано вычисление.

\newpage
\subsection{Пример: изменение влажности}

\textbf{Задача.} Относительная влажность воздуха увеличилась с 20\%
до 40\% при неизменной температуре. Во сколько раз увеличилась при этом
концентрация молекул воды в воздухе?

\medskip
\noindent\textit{Дано:}
$\varphi_1=20\%=0{,}20$, $\varphi_2=40\%=0{,}40$, $T_1=T_2=T$.

\noindent\textit{Найти:} $n_2/n_1$.

\medskip
\noindent\textit{Решение.} Рассматриваем водяной пар в воздухе как
идеальный газ. Концентрация $n$ — число молекул в единице объёма
(м$^{-3}$); индексы 1 и 2 относятся к начальному и конечному состояниям.
Искомое отношение показывает, во сколько раз изменилась эта величина.

По определению относительной влажности
\[
\varphi=\frac{p}{p_{\text{нас}}(T)},
\]
где $p$ — парциальное давление водяного пара (Па), а
$p_{\text{нас}}(T)$ — давление насыщенного водяного пара при той же
температуре (Па). Здесь $\varphi$ выражена долей единицы.
При неизменной температуре знаменатель одинаков для обоих состояний.

Чтобы перейти от давления к концентрации, воспользуемся уравнением
состояния идеального газа в записи
\[
p=nkT,
\]
где $k$ — постоянная Больцмана (Дж/К), $T$ — абсолютная температура (К).
Подставляя эту зависимость в определение влажности, получаем
\[
\varphi_1=\frac{n_1kT}{p_{\text{нас}}(T)},
\qquad
\varphi_2=\frac{n_2kT}{p_{\text{нас}}(T)}.
\]
Разделим второе равенство на первое. Общие множители сокращаются,
поэтому
\[
\frac{\varphi_2}{\varphi_1}=\frac{n_2}{n_1},
\qquad
\frac{n_2}{n_1}=\frac{\varphi_2}{\varphi_1}.
\]
Искомое отношение выражено через данные задачи. Обе части безразмерны:
единицы концентрации в отношении $n_2/n_1$ сокращаются, а относительная
влажность сама является отношением давлений. Теперь можно подставить числа:
\[
\frac{n_2}{n_1}=\frac{0{,}40}{0{,}20}=2.
\]
\noindent\textbf{Ответ:} концентрация молекул воды увеличилась в 2 раза.

На рисунке показан смысл результата: в одинаковом объёме воздуха
молекул водяного пара стало вдвое больше. Точки изображают только
молекулы воды, остальные молекулы воздуха не показаны.

% !TeX root = fundamentals.tex
\begin{center}
\begin{tikzpicture}[book diagram,x=1cm,y=1cm,
  every node/.style={font=\fontsize{10}{12}\selectfont,inner sep=2pt}]
  \foreach \origin in {0,4} {
    \draw (\origin,0) rectangle ++(2.2,1.35);
    \node at (\origin+1.1,1.65) {$V,\ T=\mathrm{const}$};
  }
  \foreach \x in {.4,1.1,1.8} {
    \foreach \y in {.38,.96} {\fill[bookmotion] (\x,\y) circle[radius=1.4bp];}
  }
  \foreach \x in {4.3,4.83,5.37,5.9} {
    \foreach \y in {.25,.68,1.1} {\fill[bookmotion] (\x,\y) circle[radius=1.4bp];}
  }
  \node at (1.1,-.35) {$\varphi_1=20\%,\quad n_1$};
  \node at (5.1,-.35) {$\varphi_2=40\%,\quad n_2=2n_1$};
  \draw[book arrow] (2.6,.68)--(3.6,.68);
\end{tikzpicture}

{\small Рис. 0.1. Сравнение равных объёмов воздуха. Число точек условно.}
\end{center}


\newpage
\subsection{Что позволяет проверить решение}

Полученное выражение полезно проверить ещё до вычислений.
В примере обе части равенства $n_2/n_1=\varphi_2/\varphi_1$
безразмерны. Это согласуется с тем, что мы ищем отношение концентраций.
Если бы после преобразований в ответе появилась единица давления
или температуры, такая запись требовала бы исправления.

Согласованность единиц необходима, но сама по себе ещё не доказывает
правильность решения. Важно проверить и физический смысл:
при постоянной температуре увеличение влажности означает увеличение
концентрации водяного пара. Наш результат больше единицы и этому
не противоречит. Если бы влажности оказались равны, полученная
зависимость дала бы $n_2/n_1=1$ — концентрация не изменилась бы.

Сравним с записью «$40:20=2$» или «пусть $x$ — ответ,
$40=20x$, значит $x=2$». Число получается тем же, но не объяснено,
почему отношение влажностей равно отношению концентраций.
Именно эту связь нужно было установить с помощью определений
и уравнения состояния. Условие постоянства температуры здесь
существенно: при разных температурах прежнее сокращение общих
множителей было бы невозможно.

Полная запись решения позволяет восстановить рассуждение:
какую систему рассматривали, какие зависимости использовали,
как получили искомую величину и что означает ответ. Пояснения
нужны в местах выбора закона и перехода к новому выводу;
каждый арифметический шаг отдельно описывать не требуется.

\medskip
{\small\itshape Условие примера: открытый банк ФИПИ, задача DDF044.
\href{https://ege.fipi.ru/bank/questions.php?search=1\&proj=BA1F39653304A5B041B656915DC36B38\&qid=DDF044}{Источник}.
Текст параграфа, рисунок и решение составлены GPT; это учебный разбор,
а не официальные рекомендации или решение ФИПИ.}
